\section{一项轮级信用受控研究的经验（ToolCredit）}\label{sec:toolcredit}

信用分配（credit assignment）方法通常由作者自己评测，对比的基线调参更少，rollout 算力也往往更少。本节报告在我们的受控研究 ToolCredit 中，用一套刻意从严的协议检验轮级信用的结果。多数结果是阴性结果（null result），我们也按阴性结果如实报告。

\paragraph{实验设置}
任务是竞赛风格的数学题，借助沙箱化的 Python 解释器（禁用网络、限制资源）求解，每个回合最多四次工具调用、预算 3{,}072 个 token，运行在 veRL 的多轮智能体循环中 \citep{sheng2025hybridflow}。策略为 Qwen3-1.7B \citep{yang2025qwen3}，先用 LoRA \citep{hu2022lora} SFT 冷启动，数据是从 Qwen3-8B 蒸馏、经拒绝采样得到的 6{,}056 条工具使用轨迹（原始模型拿到工具后反而下降 13.95\,pt）。训练使用 5{,}403 道通过污染检查的 MATH Level~3--5 题目，在单张 GH200 GPU 上进行。基线 E3 是轨迹级 GRPO \citep{shao2024deepseekmath}（64 个 prompt $\times$ 8 个 rollout，200 步，奖励为答案正确加格式正确），它把同一个结果优势（advantage）广播给每一轮。所有实验臂都从同一个 checkpoint 出发，且只改动一个组件（通过比对解析后的配置确认）。终点指标是固定集合 FULL760（MATH500、200 道 GSM8K、AIME 2024/2025）上的贪心 pass@1，并给出配对、按来源分层的 bootstrap 95\% 置信区间（10k 次重采样）。每个实验臂的判读规则（下文的情形 A/B/C）都在评测前预注册（pre-registered）；每个实验臂只跑一个 seed。

\begin{table}[t]
\centering
\small
\caption{ToolCredit 研究的各实验臂。准确率为 FULL760 上的贪心 pass@1，其后是相对 E3 的配对 $\Delta$（单位：点）及其 95\% 置信区间。“F100”表示该臂未跑 FULL760，结果取自训练期间的 100 题 MATH500 监测集。“4-call”指 FULL760 轨迹中用满四次工具调用的比例。}
\label{tab:toolcredit}
\setlength{\tabcolsep}{3pt}
\renewcommand{\arraystretch}{1.1}
\begin{tabularx}{\linewidth}{@{}l>{\raggedright\arraybackslash}X>{\raggedright\arraybackslash}p{2.2cm}>{\raggedright\arraybackslash}X>{\raggedright\arraybackslash}X@{}}
\toprule
实验臂 & 改动 & 准确率 & 关键行为效应 & 结论 \\
\midrule
E3 & 轨迹级 GRPO & .733 & 平均 0.89 次调用；4-call 1.3\% & 参照 \\
E4-A & + 执行奖励 & F100 .77（E3 .76） & 调用增多，出现疑似奖励破解样本 & 无稳定增益；奖励被利用 \\
E4-B & + 执行奖励与预算惩罚 & F100 .76 & 惩罚只能微弱抑制冗余调用 & 无稳定增益 \\
E5-v1 & $A_t{=}A_{\text{traj}}{+}0.5(s_t{-}\bar s_t)$，非守恒 & .722；$-1.05$ [$-3.68$, $+1.58$] & 4-call 92.9\%；截断率 100\% & 早期更快，终点无增益；过度调用工具 \\
E5-v2 & 零和、按 token 加权、零方差门控 & .741；$+0.79$ [$-1.84$, $+3.42$] & 4-call 0.8\%；平均 0.60 次调用 & 情形 A：未检测到增益 \\
E6 & 不屏蔽工具输出 token（80 步） & F100：未崩溃 & 4.66\% 的损失 token 来自工具返回 & 无灾难性退化 \\
E7-eq & DAPO 过滤，第 95 步（E3 的 rollout 预算） & .729；$-0.39$ [$-2.89$, $+2.11$] & 平均 0.63 次调用 & 情形 C：按 rollout 计无增益 \\
E7-200 & 同上，第 200 步（$2.53\times$ rollout） & .784；$+5.13$ [$+2.63$, $+7.63$] & 增益集中在 MATH L5、AIME & 不能归因于过滤 \\
E3-NoTool & 去掉工具 & .724；$-0.92$ [$-3.55$, $+1.71$] & --- & 情形 B：工具不带来准确率提升 \\
\bottomrule
\end{tabularx}
\end{table}

\paragraph{经验一：信用设计就是奖励设计}
在 E5-v1 中，$s_t{=}1$ 表示该次工具调用执行成功且其结果在后文被采用。这一奖励项以同一轮位置上的组均值为中心，因此在单条轨迹内并不求和为零，实际上相当于一个隐式的稠密奖励。到第 200 步，每条轨迹的净修正量随调用次数增加而上升（不调用时为 $-0.41$，调用四次时为 $+0.01$）；在零方差组（占全部组的 48\%）中，它是唯一的梯度来源。策略恰好学会了这个信号所奖励的行为：4-call 比例从 1.3\% 升到 92.9\%（采样时为 97.2\%），重复代码比例从 2.9\% 升到 92.4\%。在 F100 上早期学习更快（第 0--100 步的 AUC +3\,pt），但终点没有增益。E4 的显式执行奖励也以同样方式被利用。非守恒的信用信号本质上就是奖励，优化器会像钻任何奖励的空子一样钻它的空子。

\paragraph{经验二：守恒消除了病态行为，但短 horizon 下可改进的空间不大}
E5-v2 只改动修正项：它以轨迹自身按 token 加权的均值为中心，因此在每条轨迹内求和为零（残差 ${\le}1.1{\times}10^{-13}$），并且在结果方差为零的组中不施加。过度调用工具（over-calling）随之消失（截断率 100\%$\to$8.7\%），说明 v1 的行为漂移和早期加速都来自隐式奖励，而非轮级信用本身。然而准确率与 E3 无法区分（早期 AUC 增益 $+1.87$\,pt，低于预注册的 2\,pt 阈值）。重分配只作用于至少含两次可区分调用的轨迹；在四次调用的预算下，策略收敛到只调用一次或不调用（48.6\% 的评测轨迹没有调用），而受到任何修正的轨迹所占比例（\emph{暴露度}，exposure）从 10.4\% 降到 4.2\%。证据支持的结论是“在这一 horizon 和暴露度下未检测到增益”，而不是“信用分配没有用”。

\paragraph{经验三：在等 rollout 算力下比较}
零方差组（式~(\ref{eq:zerovar})）平均占 E3 全部组的 44\%。E7 加入 DAPO 式的动态采样 \citep{yu2025dapo}，丢弃这些组并重新采样。训练过程中零方差组比例从 34.5\% 升至 61.4\%，每步所需批次从 2.0 升至 3.06，即策略越好，过滤越贵。在 E3 的 rollout 预算下（第 95 步），终点持平；以累计 rollout 为横轴，其学习曲线落后于 E3（AUC $-1.72$\,pt）。第 200 步的 $+5.13$\,pt 增益花费了 $2.53\times$ 的 rollout，且由于缺少同预算的 E3 臂，不能归因于过滤。每次更新消耗更多 rollout 的方法，应当同时以更新步数和 rollout 数为横轴作图。

\paragraph{经验四：策略越强，工具价值越小}
E3-NoTool 臂把从 SFT 到 E3 的 $+12.76$\,pt 增益拆成三部分：6.05\,pt 来自不再被工具拖累，4.74\,pt 是真正的推理提升（置信区间 $[+2.11, +7.50]$），只有 1.97\,pt 是工具的边际价值。工具的净效应（同一组权重下有工具减无工具）从 $-13.95$（原始模型）到 $-6.05$（SFT）再到 $+1.97$\,pt（E3，置信区间 $[-0.79, +4.74]$），始终没有明确为正，且只在 MATH Level~4/5 上有帮助（$+7.03$/$+2.99$\,pt）。当单次调用可有可无、又很少起决定作用时，值得分配信用的轮本就不多，这使得短 horizon 的工具辅助数学成为检验信用方法的薄弱试验台。

\paragraph{与外部工作的关系}
这些经验与 \citet{zhang2026credit} 的重放审计及七臂剂量研究（第~\ref{sec:tax-failures}~节）一致：后者发现不同信用规则之间的差异主要反映训练剂量\unverified。这一致性佐证了我们的经验，但也削弱了其新颖性：剂量混杂已经公开。经验三同样适用于 Memory-R2：就我们能核实的范围，其 LoGo-GRPO 并未与获得相同总 rollout 数的 GRPO 比较（第~\ref{sec:memory-critique}~节）\unverified。ToolCredit 仍有三点独特之处：场景是数学加解释器，而非具身文本游戏；记录了过度调用工具的病态行为及其守恒修复；以及一套等算力协议，包含更新与 rollout 双横轴、冻结的评测和预注册的判读。要发展成论文，它还需要 horizon 扫描、一个前期调用会改变后续调用可做之事的有状态环境（这样暴露度不会衰减）、更多 seed，以及缺失的同预算基线。这些要求决定了第~\ref{sec:agenda}~节的议程。
